\title{Distributed Shared Timestamped Memory}

%\author{\IEEEauthorblockN{Anonymous Author(s)}
%\IEEEauthorblockA{\textit{Affiliation Withheld for Double-Blind Review}}}

\author{
\IEEEauthorblockN{Nicholas Cassarino, Swann Perarnau, Jonathan Ozik,\\
Nicholson Collier, Jamie Cook, Kevin Brown}
\IEEEauthorblockA{\textit{Argonne National Laboratory} \\
Lemont, IL, USA \\
\{ncassarino, swann, jozik, ncollier, jcook, kabrown\}@anl.gov}
}

\maketitle

\thispagestyle{fancy} 
\lhead{} 
\rhead{} 
\chead{} 
\lfoot{\footnotesize{ 
SC26 Workshops, November 15-20, 2026, Chicago, Illinois, USA 
\newline 979-8-3195-1221-5/26/\$31.00 \copyright 2026 IEEE}} 
\rfoot{} 
\cfoot{} 
\renewcommand{\headrulewidth}{0pt} 
\renewcommand{\footrulewidth}{0pt} 

\begin{abstract}

% 150 word maximum
%   this below is 149 words
Parallel discrete-event simulation (PDES) and agent-based modeling and simulation (ABMS) are two important modeling and simulation paradigms used in domains such as network simulation, 
epidemiology, and transportation. Both must coordinate distributed state as simulation time
advances, yet common approaches in each paradigm sacrifice
either performance through blocking synchronization or scalability and simplicity
through rollback mechanisms to do so. Moreover research in each field has largely progressed in isolation.

We present a partitioned distributed versioned array that can serve as a
common global distributed memory layer in both PDES and ABMS for more flexible synchronization
schemes. By allowing each rank to maintain multiple versions of a shared
data structure available at the same time, our approach enables unbalanced
workloads to drift out of lockstep without rollbacks or neighborhood barriers.
We demonstrated increased performance using this scheme on common benchmarks
from both PDES and ABMS paradigms, against state of the art runtimes.

\end{abstract}

\begin{IEEEkeywords}
distributed simulation, programming models, MPI alternatives, HPC
\end{IEEEkeywords}

\section{Introduction}          % problem, motivation, contributions

% Intro: PDES and ABMS
%        Looking provide a programming model that can provide High-Performance
%             on these two growing simulation fields
Parallel discrete-event simulation (PDES) and agent-based modeling
and simulation (ABMS) are two important modeling and simulation paradigms 
used to study physical systems in domains such as network simulation~\cite{network-sim}, epidemiology~\cite{epi}, and transportation~\cite{transport}.
The two paradigms target different systems: PDES
handles systems with well-defined states and processes that drive
state transitions between them, while ABMS handles systems of
autonomous agents whose local behaviors and interactions
collectively evolve the global system state. These differences in focus have
caused each paradigm to scientifically progress separately, developing
their own runtimes, programming models, and scalability techniques
rather than leveraging each other's successes.

% Problem: can look at the proposal for this
%   problem is that existing techniques use either Conservative or Optimistic
%       the conservative case ensures correctness but suffers performance wise
%       the optimistic case can struggle with scaling due to cascading rollbacks
%   instead of providing a hybrid model, we provide a distributed programming model 
%       that can be used for both paradigms effectively. Therefore, this section should
%       be framed differently, we do not need to build up to a hybrid model but instead 
%       focus on the fact that there exists no good model that can be used by 
%       both communities and avoids siloing

% most visible in how each paradigm handles a distributed
%   setup...
% We see this siloing in how each paradigm handles distributed simulation
%   suites
Siloing can be observed in how each paradigm approaches distributed
simulation. Conservative PDES guarantees correctness
by blocking each process until it can prove no earlier
event will arrive, which ensures correctness but forces frequent synchronization.
Optimistic PDES
eliminates blocking by processing events eagerly and rolling
back on causality violations. At scale optimistic PDES suffers from cascading
rollbacks, state-history overhead, and the programmer burden of
writing reverse handlers. 
ABMS frameworks have approached distributed simulation differently. A notable approach involves 
advancing all ranks in lock-step through a per-tick ghost exchange
that serializes boundary agents and neighbor-side state, transmits
them via a blocking point-to-point sends to each neighbor, and deserializes them as read-only
ghost copies before the next tick begins. This gives agents
uniform access to neighborhood state at the cost of frequent synchronization,
with further optimization
left to the application.

% hmmm not sure how i feel about this section
%   i like that it brings up the gap that we will fill
%   not sure of it's effectiveness

Each of these approaches was developed in response to the needs of its own
community, and none generalize cleanly to the other. The result is that a
scalability advance in one paradigm, whether a smarter conservative window for
PDES or a more efficient ghost exchange for ABMS, rarely benefits the other.
%
We argue here that some scalability techniques can benefit both communities by
expressing them as distributed data structures with well tailored synchronization
mechanisms. While this approach might not provide a distributed
programming model general enough to serve both paradigms yet, it offers an avenue
for the gradual design of a common substrate so that innovations can be expressed once
and reused across communities.

%What is
%missing is a distributed programming model general enough to
%serve both paradigms on a common substrate, so that scalability
%techniques can be expressed once and reused across communities.

% Motivation: DSTM 
%   we believe that DSTM can provide a unique distributed programming model that is 
%       efficient and usable for both PDES and ABMS
%   the space-time memory model (versions) provides a unqiue abstraction into
%       async distributed communications that limits synchronization (allowing 
%       for drift between ranks). the distance matrix abstraction contains
%       allows for stronger control over ranks synch/communication. We 
%       discuss ideal scenarios for performance with this model. 
%   the space-time memory model provides a conservative simulation technique
%       that loosely bounds ranks by 'distance' and local neighborhoods. this
%       allows for low computation ranks to speed ahead. low sync allows for
%       high performance in both PDES and ABMS

We argue that a distributed programming model organized around
\emph{timestamped memory}, one that versions distributed state and
differentiates for distributed operations the virtual simulation time at
which they occur from the virtual simulation time they become visible,
can serve both PDES and ABMS paradigms effectively. These timestamps then
give the opportunity for a runtime to preserve causality (ensure a version is not read until no operation that could modify it can still arrive) without either
global barriers or rollback, when associated with a version history that
lets producers and consumers reference the same state at different points
in simulation time.

We demonstrate in this paper a realization of this idea in DSTM, a distributed
service that provides a partitioned distributed array data structure built on three
abstractions: as each rank in a simulation owns a partition of this array, we
enable (1) \emph{versioning}, in which each rank retains several versions (i.e.
a history) of its state tagged over simulation time; (2) \emph{timestamped
operations}, in which reads and writes to the distributed array are tagged with
the simulation time at which they occur; and (3) a \emph{distance matrix} that
describes which nodes have the potential to send operations between each other
and how long an operation would take (in simulated time) to become visible.

Together these abstractions let ranks drift out of lockstep with no rollbacks
and no global barriers: lightly-loaded ranks race ahead and deposit future
versions, while lagging ranks read the earlier versions their producers have
already made visible. Bounds on reachable neighbors preserve correctness. The
result is an execution model whose synchronization cost is bounded by local
simulation-specific neighborhood structure and per-edge drift. Making DSTM well
suited for simulations from both PDES or ABMS when both the sparse event traffic
of PDES and the tick-driven ghost exchange of ABMS.

% Contribution:
%   we present DSTM, a distributed space time memory model that provides conservative
%       high performance for both PDES and ABMS
%   we demonstrate the affect of application and distance matrix on drift size
%           "we demonstrate how application conditions effect the abstrations
%            that DSTM provides"
%   we compare this framework directly against prior works {PDES: ROSS, ABMS: Repast}

This paper makes the following contributions:
\begin{itemize}
  \item We design and implement DSTM, a distributed service
        for PDES and ABMS runtimes. Through a distributed array abstracting over versioned memory, timestamped messages,
        and a distance matrix, DSTM eliminates rollback and global barriers while
        tolerating inter-rank drift.
  \item We characterize how our distributed array data structure
        can map to a simulation structure and identify the workload regimes in which
        the approach is expected to perform best. %\hl{something to discuss in DSTM Overview or Discussion}
        
  \item We evaluate DSTM against representative prior systems
        from each paradigm, ROSS~\cite{ross, ross-recent} for
        PDES and Repast~HPC~\cite{repast-hpc} for ABMS, using
        benchmarks that reproduce identical simulation semantics
        across all three runtimes.
\end{itemize}

\section{Background and Related Work}
\label{sec:background}

\subsection{Simulation Paradigms}

The two large-scale simulation paradigms at the focus of our
work, parallel discrete-event simulation (PDES) and agent-based
modeling and simulation (ABMS), are both commonly deployed in
distributed settings but express computation in fundamentally
different ways.

\emph{PDES} advances a global virtual clock by processing
timestamped events across logical processes (LPs). Each LP owns
a piece of state that changes only in response to events it
receives, and the runtime is responsible for preserving causality
across LPs. This has historically been done through one of two
families. \emph{Conservative} PDES, originating from Chandy et
al.~\cite{chandy-misra}, only processes an event once every
neighbor has advanced far enough to ensure it is safe. Safety is
established through per-link lookahead and distributed via null
messages, ensuring correctness but requiring frequent synchronization. 
\emph{Optimistic} PDES, introduced as Time
Warp~\cite{Time-Warp}, instead processes events eagerly and
recovers from causality violations through rollback and
antimessages, tolerating unbounded rank drift but paying for it
with rollback overhead, state history for recovery, and the user
burden of writing reverse handlers. ROSS~\cite{ross, ross-recent}
supports both strategies and serves as our PDES baseline.

\emph{ABMS} instead distributes stateful autonomous agents
across ranks and advances them in lock-step through discrete
ticks. Agent behavior is typically local, but often depends on
neighborhood or global state exposed through shared spatial
projections such as grids, networks, or continuous spaces. Unlike
PDES, where progress is measured in events and causality is
enforced per event, ABMS typically advances \emph{all} ranks
together. Each tick, every agent executes its behavior, and a blocking neighbor exchange gathers ghosted state from partition-adjacent ranks before the next tick begins. Repast HPC~\cite{repast-hpc}
and its Python successor Repast4Py~\cite{repast4py} provide this
pattern, as a default, with an explicit per-tick \texttt{synchronize()} that
exchanges agent and projection state across partition boundaries. Further optimizations
to communication and synchronization are left to the user. 

\subsection{Space-Time and Time Warp}

DSTM's use of timestamped state has roots in the space-time view of
distributed simulation. Jefferson's \emph{virtual time} and its
realization in Time Warp~\cite{Time-Warp} framed a distributed
computation as points on a two-dimensional space-time plane, with
each process advancing along its own local virtual clock and rolling
back when a straggler message revealed a causality violation.
Fujimoto's analysis of Time Warp under synthetic
workloads~\cite{fujimoto1990phold} characterized how rollback
frequency, state-saving cost, and rank drift interact to determine
throughput, and established that tolerating drift is only
worthwhile when the mechanism for recovering from it is cheap.

More recent systems have revisited the space-time abstraction
without committing to full optimistic rollback. Chandy and
Sherman's space-time simulation~\cite{chandy-sherman} treated
the simulation state as a two-dimensional grid indexed by process
and time, computed by relaxation rather than event scheduling.
DSTM inherits this grid-shaped view of state (one version per
(rank, time) pair) but replaces speculative execution and
rollback with distance-bounded visibility, so drift is tolerated
without ever having to undo work.

\subsection{Computing with Drift}

% A related concept to DSTM's multi-version drift comes from
%   transactional systems. (MVCC) produces a new version tagged with a
%   timestamp on every write rather than overwriting in place. 
% This allows readers to see a consistent snapshot at their own
%   timestamp without blocking concurrent writers. MVCC allows for
%   drift between readers and writers but pays 
A related concept to DSTM's versioned drift comes from transactional
systems. 
%The most closely related concept to DSTM's versioned state comes
%from outside HPC entirely. 
Multi-Version Concurrency Control
(MVCC)~\cite{bernstein-mvcc} produces a new version tagged with a commit
timestamp on every write rather than overwriting in place. This
allows readers to see a consistent snapshot at their own
timestamp without blocking concurrent writers. MVCC allows drift
between readers and writers but pays for it with conflict
detection and transaction abort at commit time. %, semantics a
%tightly-coupled simulation cannot afford.

DSTM builds upon the versioned-state idea and applies it to
distributed simulation by tying versions to simulation timesteps
rather than commit timestamps, and by bounding visibility through
the distance matrix. This avoids the abort semantics MVCC
relies on while preserving drift tolerance.

\subsection{Asynchronous Runtimes and PGAS}

A parallel line of work has attacked the same MPI+X limitations
from the systems side rather than the simulation side.
UPC++~\cite{upcxx} provides one-sided \texttt{put}/\texttt{get}
over a global address space, with futures for asynchrony but no
notion of logical time. Charm++~\cite{charmpp, charm14} uses
migratable message-driven \emph{chares} whose asynchronous
execution tolerates load imbalance well and naturally scopes
communication to neighborhoods. It has served as the substrate
for simulation codes such as ChaNGa~\cite{changa}, but exposes
no versioned state and leaves ordering to message handlers.
HPX~\cite{hpx-lci} and Lamellar~\cite{lamellar} are asynchronous
many-task runtimes paired with a PGAS-style shared address space.
None of these systems expose per-timestep versioning or per-edge
visibility bounds. All leave temporal semantics to the
application.

DSTM shares the one-sided \texttt{put}/\texttt{get} interface of
these systems but adds versioned state and distance matrix bounded
visibility as abstractions, so time-stamped delivery
is provided by the runtime rather than reimplemented in every
application.

\section{DSTM Overview}         % programming model, key abstractions, API

% We have room here to described our three abstractions
%   Versioned memory
%   Timestamped messages
%   A distance matrix

While we envision DSTM as an approach that can eventually become a full-fledged
programming model and runtime, we focus our work here on providing a
partitioned, distributed, versioned array as a data structure that can be
effectively kept synchronized across both PDES and ABMS simulations.

As an initial step towards a more general approach, we make several simplifying
design choices: a simulation is expected to be distributed among a set of static
and uniquely identifiable ranks, with a unique shared virtual timeline across
ranks (while each rank progresses along the timeline at its own pace, it is the
same simulation timeline). Virtualization provided by logical processes, or
virtual agents, is left for future work. Similarly, simulation state is assumed
to be partitioned among ranks statically (no load balancing), while each rank
only exhibits a single simulation time at once (no internal virtualization of
the timeline for logical process management for example).

\begin{figure}[t]
    \centering
    \includegraphics[width=\linewidth]{figs/darray_final.drawio.png}
    \caption{Per-index versioning across ranks. Each rank owns a contiguous slice of the global index space (rows) and maintains an independent version per timestep (columns). Cells are colored by state.}
    \label{fig:DSTM-darray}
\end{figure}

\subsection{Core Abstractions}

From these assumptions, we build our core abstractions. First, a \texttt{darray}
is provided as a 1D array of identically sized elements distributed across all
simulation ranks. Distribution allows for an uneven contiguous slice of indices
on each rank. This distributed array is backed by \textit{versioned memory}: any
operation to an index is associated with a timestamp, and multiple
versions of each index can exist simultaneously.

Second, distributed operations are provided to access this distributed array.
These operations are \texttt{put} (write) and \texttt{get} (read). 
%While they target a specific index the source, respectively destination, of the data written, respectively read, is a local buffer.
Each operation targets a specific index in the distributed array. For a put, the source of the data is a local buffer; for a get, the destination is a local buffer.
Operations are also augmented with two timestamps: their originating time (the simulation time of the originating rank) and their destination time or delta (the time the operation takes to become visible). 
Destination time has a subtle meaning: for writes, they are a distance forward in time that it will take for the operation to appear; for reads, they are a distance backward in time, an intent to read from an older state. 
In both cases, local operations are allowed up to the current version. \texttt{Put} operations must be in the future (or present for local operations), as a \texttt{put} into the past may have already been passed by a neighboring (receiving) rank. 
In the same way, \texttt{get} operations must be in the past (or present for local operations), as a \texttt{get} into the future may cause catastrophic blocking. 
The opposite of these cases (\texttt{put} in past; \texttt{get} in future) do not align with our current paradigm but do spur interesting discussion.

Finally, we also assume that the simulation structure (e.g. the network
topology, geographic distribution of airports, transport hubs) implies (1) a
sparse communication structure between ranks, that is, across the full runtime
of a simulation, not every rank will require communicating with every other
rank; (2) a distance between ranks: a number of simulation timesteps that an
action will take to become visible to another rank. These last two assumptions
are not limiting: a full communication graph and a distance of 0 between
all ranks is a correct setup for our approach (equivalent to all-to-all
synchronization at each timestep), at the cost of nullifying the optimizations it allows.
This communication graph is communicated to the DSTM as a \texttt{distance matrix} storing the distance between each rank pair, with 0 for the distance between local ranks and null for ranks not reachable from each other. Distances are assumed symmetrical. Using the distances within our communication graph, we are able to restrict the progress on quick ranks by the distance between its neighbors (hence a full communication graph of distance 0 being equivalent to all-to-all).

%\hl{to clarify, using the distances within our communication graph we are able to restrict the progress on quick ranks by the distance between it's neighbors (hence a full communication graph of distance 0 being equivalent to all-to-all). another clarification, \texttt{put} operations must be in the future (or present for local operations), as a \texttt{put} into the past may have already been passed by a neighboring (receiving rank). in the  same way \texttt{get} operations must be in the past (or present for local operations), as a \texttt{get} into the future may cause catastrophic blocking. the opposite of these cases (\texttt{put} in past; \texttt{get} in future) align with our current paradigm but do spur interesting discussion.}

While we describe later in this document specific use cases, we can identify two
main synchronization patterns using these abstractions: a pull model and a push
model. In the pull model, each rank submits updates (puts) to local data at the
current version, while neighbors pull data by issuing gets in the past. In the
push model, each rank pushes updates to its neighbors in the future, while
retrieving (gets) local information at the current version. In both
cases, the scheme enables ranks to drift slightly apart.

% We can use this figure below as a reference to how some of these features work
%   We have a figure showing maximum drift in a simulation
%   We have a figure showing total active versions per rank
%%      Both of these on a ring model:
%           The ring model could be a simple way to demonstrate what we are doing
%           And we have the diagram above.
%%%%%%% If you would like to reference any of that data while writing this it 
%           resides within examples/minimal/ring_mini/profiling
%           &              variable_mini has some ones on weight graph topographies

\subsection{Version Availability}

Beyond the specific abstractions for this distributed array data structure, the
DSTM can be understood as an asynchronous messaging layer on top of which
operations to arrays are issued. The DSTM is informed on the local time of each
rank by two mechanisms: (1) a rank is only allowed to issue operations with a
monotonically increasing originating time; (2) heartbeats are sent to all neighbors
after each increase in local time. From this information about the progress in
time of a simulation, rules about when operations are allowed to complete can be
derived. We identify for each version of an index the following states:
writable, readable, discardable.

\paragraph{Writes} a version $t$ of an array index is created by any write of
destination time $t$. It stays writable (available for further writes) until no
neighbor, nor the owning rank, can issue a put to it (progressed in originating
time past this version at time greater than $t - dist$). As a simplification,
overlapping writes are ignored (not a conflict).

\paragraph{Reads} a version $t$ of an array index become readable when it is no
longer writable, and stay readable until no neighbor, nor the owning rank, can
issue a get to it (progressed in originating time past this version at time
greater than $t + dist$). At which point it become discardable.

These visibility rules have an intuitive meaning: they use the distance between
two ranks as an upper (or lower) limit to how long an operation can take to
reach the destination rank. For example, when simulating a network, the time a
message takes to transmit between routers limits how often two routers would
need to synchronize.

We also note that discarding versions is a memory management optimization: it
allows the DSTM to maintain less versions of each array index, without changing
the semantics of the program.

% We can also use a figure or two to visualize how drift changes with 
%   application structure/distance matrix (neighborhood)
%       This is bullet point two of the contributions
%   we can also talk about how many versions are active depending on 
%       neighborhood size.

\begin{figure*}[t]
    \centering
    \includegraphics[width=0.85\linewidth]{figs/big_diag.drawio.png}
    \caption{An example execution graph of a simulation with a ring topology ($n=4$)}
    \label{fig:DSTM-diagram}
\end{figure*}

\section{Implementation}        % runtime, comm layer, data movement

% This section could be combined with the prior, if not would contain:
%   low-level details about communication and barrier implementations
%       3 barrier

The DSTM is implemented as a prototype messaging layer on top of MPI, with
individual messages between ranks (operations, barriers, and heartbeats) as
asynchronous point to point operations containing a message header with
timestamps and distances. The current implementation makes no attempt to improve
on the granularity of operations, generating 1 message for puts and 2 messages
for gets (request-reply).

We point out here two specific implementation details of the DSTM as critical for its operation:
\begin{itemize}
    \item minimum virtual time: the DSTM maintains an internal virtual time clock that only advances to match the smallest
    version available (not discarded) across all its own partitions.
    \item barrier/termination detection: a logical ring algorithm is used to identify when the simulation finishes.
\end{itemize}

The barrier ring algorithm merits more inspection: we have indicated that local
time information is extracted by the DSTM from the originating time of locally
issued operations. In order to detect the simulation is done, one would have to
progress past the last timestamp of any operation. And any pending operation
would have to complete by detecting that neighbors have progressed past that
last timestamp. Thus, we implement a 3 phase barrier algorithm instead of the
more classic 2 phases to identify the case where all ranks have \textit{sent out} all
their operations from the case where all ranks have \textit{completed} their operations.
Thus, all ranks arranged in a logical ring, the first round the token is
forwarded when all local operations have been issued. The second round around
the ring, the token is forwarded when all local operations have completed, and
the third indicates exit from the barrier. These barriers tokens are also taken into account when updating the visibility of a version.

The model in brief
An address in DSTM is (rank, offset, time). Time is logical simulation time (dstm_time_t), not wall clock.
Every operation carries a (time, delta) pair; the version it touches is time + delta.
A put requires delta >= 0 (you cannot write the past), a get requires delta <= 0 (you cannot read the future). delta == 0 is legal only for the local rank; remote accesses must use a non-zero delta.
Versions do not inherit from one another. A cell nobody wrote at time t reads as zero, even if it was written at t-1. Long-lived state has to be re-published every tick. This is the most common first-use mistake.
A dstm_distance_t is the communication graph plus a per-edge lookahead: how many timesteps of slack exist between two ranks. It is what lets a rank decide a version is final without asking anyone. A lookahead of 0 is unsound; use 1 for "synchronise every tick".
put is asynchronous. get blocks until the requested version is readable, and that block is the only synchronisation in the system.

## Most recent accurate context:

\section{An Overview of DSTM}
\label{ref:overview}
The goal of our runtime/programming model is to express locality and sparsity in order to allow faster ranks to race ahead of slower ones. We achieve this by representing time as a dimension with our data structures, allowing for communication to occur at a specific timestamp. In this case, each rank has a history of data over time, allowing for writes to be asynchronously in the future and reads to not blocking the rank being read from. Our attempt to achieve this goal has resulted in the abstractions mentioned within this section, while implementation details follow in Section~\ref{ref:implementation}.

\subsection{Timestamped Operations}

% puts, gets, append, accumulate

The array, graph, and queue interfaces are built on four timestamped
communication operations (Table~\ref{tab:timestamped-operations}):
\texttt{get}, \texttt{put}, \texttt{append}, and \texttt{accumulate}. Each operation specifies an originating
simulation time $t$ and a time offset $\Delta$, producing an effective
timestamp
\[
    t_{\mathrm{eff}} = t + \Delta .
\]
The operation accesses the version of the destination region associated
with $t_{\mathrm{eff}}$, rather than the destination rank's current
simulation time.

A \texttt{put} writes state at the effective timestamp and completes
asynchronously. Puts may target the future ($\Delta > 0$), or the local
present ($\Delta = 0$ for a local operation). A remote put cannot target
the present, because the destination rank may have already made that
version readable. Similarly, puts cannot target the past.

A \texttt{get} reads state at the effective timestamp and blocks until
that version becomes readable. Gets may target the past
($\Delta < 0$), or the local present ($\Delta = 0$ for a local
operation). Remote gets cannot target the present or future because the
requested version may not yet be complete.

%paragraph for accumulate/append since each data structure can use either
%   accumulate or append
In addition to overwriting state, we extend DSTM by providing \texttt{append} and
\texttt{accumulate} operations for combining multiple writes to the same
effective timestamp. The queue uses \texttt{append} to concatenate
contributions from multiple producers; their order is determined by message
arrival and is therefore not guaranteed. Arrays and graph vertices and edges
provide \texttt{accumulate}, which combines contributions using a
user-registered reduction function. The first contribution initializes the
target value, and subsequent contributions are folded into it. Both operations
follow the timestamp rules of \texttt{put}: they are asynchronous and may
target the future, or the local present. A program must not mix
\texttt{put} and \texttt{accumulate} on the same value and effective
timestamp.

\begin{table}[t]
\centering
\caption{Timestamped DSTM operations. This work extends the earlier
\texttt{put}/\texttt{get} interface with \texttt{append} and
\texttt{accumulate}.}
\label{tab:timestamped-operations}
\begin{tabular}{llll}
\toprule
Operation & Effect & Remote $\Delta$ & Completion \\
\midrule
\texttt{get}        & Read version          & $<0$ & Blocking \\
\texttt{put}        & Overwrite value       & $>0$ & Asynchronous \\
\texttt{append}     & Concatenate value     & $>0$ & Asynchronous \\
\texttt{accumulate} & Reduce into value     & $>0$ & Asynchronous \\
\bottomrule
\end{tabular}
\end{table}

\subsection{Versioned State}
\label{ref:versions}
Each DSTM data structure is backed by a region containing multiple
versions of its state. A version is identified by its simulation
timestamp and moves through three states: writable, readable, and
discardable. Puts, appends, and accumulates contribute only while a version is
writable. Gets complete after the version becomes readable, indicating
that no permitted writer can still modify it. Once no rank can issue a
valid get for the version, it becomes discardable and its storage may
be reclaimed. The same version lifecycle applies to arrays, graphs, and
queues.

\subsection{Distance Matrix and Progress}
DSTM represents communication locality using an application-provided
distance matrix. For a pair of ranks, a positive entry specifies the
minimum number of simulation timesteps separating an operation from its
effect on the destination rank. A negative entry indicates that the
ranks do not communicate, while the diagonal is zero. The matrix can
therefore be viewed as a weighted communication graph
(Fig.~\ref{fig:dist_matrix}).

The runtime combines these distances with the last known simulation time
of each neighboring rank. This information determines whether another
put or get could still target a particular version, and therefore when
that version may become readable or discardable. Larger distances allow
more drift between ranks, while shorter distances impose tighter
progress constraints.

\begin{figure}[t]
    \centering
    \includegraphics[width=0.8\linewidth]{diagrams/paper_dist_matrix.drawio (1).pdf}
    \caption{Communication graph derived from the distance matrix.}
    \label{fig:dist_matrix}
\end{figure}

\subsection{Programming Contract and Limitations}
%timestamp monotonicity, put/get delta rules, blocking reads, distance-matrix soundness, and lack of version inheritance
DSTM relies on three rules for safe progress. Gets are blocking, ranks issue
operations with monotonically increasing originating timestamps, and the
runtime sends heartbeats to neighbors when local time advances. Together with
the distance matrix, these rules determine when no more writes can target a
version, making it readable, and when no more gets can target it, making it
discardable. Remote writes must target the future ($\Delta>0$), while remote
gets must target the past ($\Delta<0$). Correctness requires the
application-provided distance matrix to conservatively describe these
communication constraints. Violating this contract may cause a version to
become readable or discardable before all valid operations have completed.

DSTM does not inherit state between versions. Therefore, version histories may
contain gaps, and applications must create each version they later read. We
have not found these gaps problematic; instead, they preserve temporal sparsity
by avoiding unnecessary versions and storage.

\section{Composable Timestamped Data Structures}
The DSTM programming model provides three timestamped data structures:
a distributed array, graph, and queue. These structures provide the same
timestamped put and get abstractions but differ in their communication
granularity, collision handling, and expected use. This section describes
these trade-offs and then explains how multiple structures can be composed
within a single simulation through DSTM regions.

\subsection{Distributed Array}
The simplest of the three structures is the distributed array.
DSTM versions each index of the array, creating a two-dimensional
space of indices and timestamps. With \texttt{put}, writes by different
ranks to the same index and timestamp may overwrite one another.
Applications can avoid these collisions by assigning a separate index
to each writer or by using \texttt{accumulate} to combine contributions
with a user-provided reduction function. However, applications must
manage this indexing or reduction explicitly.

\subsection{Distributed Graph}
The distributed graph versions both vertex and edge state. Edges provide
a natural location for communication between neighboring vertices,
reducing the need for applications to construct a multi-indexed array.
Vertex and edge state support \texttt{put}, \texttt{get}, and
\texttt{accumulate}; when multiple ranks contribute to the same state
and timestamp, \texttt{accumulate} combines their contributions using a
user-provided reduction function rather than allowing them to overwrite
one another. Separating vertex and edge state also supports different
forms of exchange, such as storing congestion at vertices while agents
move along edges. Graphs are therefore naturally suited to geographic
workloads, as demonstrated by our traffic simulation.

\subsection{Distributed Queue}
The queue data structure handles collisions by appending \texttt{puts} to the same index to a queue, converting them to a \texttt{push}. Each one of these indices is versioned therefore, to avoid any timestamp complexities, our \texttt{pop} operation dequeues the entire queue. \texttt{Puts} may come in random orders as well, so the user must be prepared to handle the queue within their own application. The queue suits itself towards a large variety of applications, as in someways it is an improved version of our array data structure. One example may be a network simulation, where packets are flying all over and landing in different ranks inboxes. In this case the user will not have to worry about handling a multi-indexed mailbox, which requires some hashing computation to ensure no collisions. In this case, the user does not need to manage a multi-indexed mailbox
or compute a mapping that prevents collisions.

\subsection{Multi-Region Composition}
With complex applications a programmer may be inspired to use more than one data structure, possibly even with different distances between ranks. DSTM provides this feature through regions. A region is a memory backing to a data structure, and in this version of DSTM it is possible to create multiple regions (and therefore data structures) within a single application. Users can take advantage of this to create more complex simulations, examples of which exist within Section~\ref{ref:traffic} \&~\ref{ref:datacenter}.

In cases where two data structures within an application have different distance matrices, DSTM has some more handling to do. Given a rank pair that has multiple different distances within multiple different distance matrices, the implementation is as follows. While in single data structure applications a single integer value is used to facilitate transitions between version states, multiple data structures with different weights may cause premature transition from readable to discardable (for more details on this transition refer to Section~\ref{ref:versions}). Therefore, DSTM combines all distance matrices, collecting the maximum and minimum distance for each rank pair. The minimum is used for the transition between writeable and readable, and the maximum is used for the transition between readable and discardable.

\section{Implementation of DSTM}
\label{ref:implementation}

\subsection{Architecture and Communication}
Request model
Progress mechanism
idk

\subsection{Progress and Memory Reclamation}
the optimization of discardable (this is present within workshop paper so maybe not interesting; likely should replace with barrier implementation since workshop reviewers were interested.)

\subsection{Multi-Region Distance Combination}
this is a intellectual contribution. novel %this is already in section 4.4 Multi-Region Composition. Consider deleting this subsection here.

\subsection{Version State Machine}
All of the implementation details of when these version states changes and why and how.
This is also in the workshop paper, we need to give details to how it's different for different data structures or different in the case of combined distance matrices

\subsection{Rocq Validation}
Due to DSTM's untraditional structure, efforts have been made to validate a simplification of this model. Using ROCQ (formerly known as COQ)~\cite{rocq-prover}, we can define the core abstractions and prove them. As a part of this effort we have also proved a minimal ring simulation\footnote{This proof can be found: \textit{absent for anonymity}}. We have been able to validate the abstractions and prove that a ring program can terminate. Version safety is real.

Validation gives us more confidence in our abstraction and helps us explain our ideas in a more formal method.
